% Chapter 5. GaugeCore and exact G4K2B / H5E8AA / OP2c chronology.
\chapter[Gauge Symmetry and Physical Hilbert Space]{Color Gauge Symmetry, Physical States, and Nontriviality}
\label{ch:gauge}

\section{The gauge action is part of the model's data}
For $N=2$ the gauge group is the \emph{actual} special unitary matrix
group $G=\mathrm{SU}(2)$. It acts on the nine matrix coordinates by color
conjugation and on the original fermionic fiber by a compatible unitary
representation. We should not infer the existence or properties of the
latter from the abstract fact that SO(3) rotates three color coordinates:
the original 48-Majorana/Fock implementation, including central phases,
has its own construction and proof lineage.

In the source, admissible \texttt{GaugeData} consists of maps
\begin{equation}\label{eq:gauge-actions}
 b:G\longrightarrow\mathrm{Iso}_{\mathbb R}(\mathcal X),\qquad
 u:G\longrightarrow\mathrm{Iso}_{\mathbb C}(\mathcal F)
\end{equation}
that respect multiplication and depend continuously on the group
parameter. They satisfy the exact color-coordinate intertwining
\begin{equation}\label{eq:boson-adjoint}
 X_i\bigl(b(g)x\bigr)=gX_i(x)g^{\dagger},
\end{equation}
where $X_i(x)=\sum_A x_i^A t_A$ is the original Hermitian color matrix.
The corresponding fermionic covariance is
\begin{equation}\label{eq:fermion-gauge-covariance}
 u(g)\,\theta_{\alpha A} v
 =\sum_B R(g)_{AB}\theta_{\alpha B}\,u(g)v,
\end{equation}
where the real adjoint coefficients $R(g)_{AB}$ are defined from the
\emph{source} trace of color conjugation, not from a substituted
three-dimensional representation. This is the actual source requirement
\texttt{fermion\_adjoint}.

\section{The physical Hilbert subspace is a closed kernel intersection}
Pullback by $b(g)$ and fiber action by $u(g)$ induce linear isometries
on the original $\mathscr H=L^2(\mathcal X;\mathcal F)$. Write
\begin{equation}\label{eq:physical-condition}
 \mathscr H_{\mathrm{phys}}
 :=\bigl\{\psi\in\mathscr H:
 \psi\circ b(g)=u(g)\psi\text{ in }L^2
 \text{ for every }g\in G\bigr\}.
\end{equation}
The equation is an equality of $L^2$ classes: it is not asserted at every
point of an arbitrary measurable representative.

\begin{theorem}[Closedness of the original physical space]
The source-defined space $\mathscr H_{\mathrm{phys}}$ is a closed complex
linear subspace of $\mathscr H$, hence a Hilbert space in the restricted
original inner product.
\end{theorem}
\begin{proof}
For each $g\in G$, let $P_g$ be the continuous measure-preserving
pullback on $\mathscr H$ and $U_g$ the continuous fiber action. The
equality constraint is $\ker(P_g-U_g)$, a closed complex subspace.
The intersection of closed subspaces over all $g$ is closed. This is
exactly the source's \texttt{physicalSpace} definition and
\texttt{physicalSpace\_closed} lemma.
\end{proof}
\formalrecord{\texttt{AlgebraData.GaugeData},
\texttt{GaugeData.bosonPullback}, \texttt{GaugeData.fiberAction},
\texttt{GaugeData.physicalSpace} and
\texttt{GaugeData.physicalSpace\_closed}, in \texttt{GaugeCore.lean}.
Atlas G-01.}

\section{The invariant smooth core and its stronger pointwise law}
Let
\begin{equation}\label{eq:invariant-core}
 \mathscr D_{\mathrm{phys}}
 :=\{\psi\in\mathscr D:\iota\psi\in\mathscr H_{\mathrm{phys}}\}.
\end{equation}
Because its members are continuous and compactly supported, the source
strengthens their $L^2$ equality to the pointwise law
\begin{equation}\label{eq:smooth-gauge-pointwise}
 \psi(b(g)x)=u(g)\psi(x),\qquad
 x\in\mathcal X,\;g\in G.
\end{equation}
There is a subtle measure-theoretic step here. The source first uses
the measure-preserving action to establish an almost-everywhere equality.
Both sides of \eqref{eq:smooth-gauge-pointwise} are continuous in $x$;
Euclidean Lebesgue volume has full support, so almost-everywhere agreement
of continuous functions implies equality everywhere. Only for \emph{this
smooth-core argument} can one harmlessly evaluate at a chosen $x$.

\begin{remark}[The quantifiers in the pointwise upgrade]
For \emph{each fixed} $g$, equality of the two $L^2$ classes gives
an almost-everywhere identity in $x$. Continuity in $x$ and full
support of Lebesgue measure then imply equality for \emph{every}
$x$ for that same $g$. Because this reasoning works for an
arbitrary $g$, the conclusion has the asserted quantifiers
$\forall g\,\forall x$; no uncountable union of exceptional null
sets is used. Conversely, an arbitrary $L^2$ class lacks this
continuity and cannot be evaluated pointwise merely because it
is physical. This is exactly the distinction between the source
\texttt{invariantCore_equivariant} conclusion and the defining
\texttt{physicalSpace} class equality.
\end{remark}

Taking norms and using the fermionic isometry gives the identity
\begin{equation}\label{eq:gauge-norm}
 \|\psi(b(g)x)\|_{\mathcal F}=\|\psi(x)\|_{\mathcal F}.
\end{equation}
This exact gauge-norm invariance will later support changes of color
direction in conditional confinement arguments. It does not mean that
arbitrary fixed slices of an $L^2$ wavefunction are defined.

\section{Is the physical Hilbert space nonzero?}
A definition as the intersection of symmetry constraints is not enough:
it could yield the zero space. We need a \emph{normalizable nonzero} state,
not just a formal gauge-invariant vector in the finite fermionic fiber.
The historical FCP source produced such a witness.

\begin{theorem}[Nonzero physical $L^2$ state;
historically qualified]\label{thm:physical-nonzero}
For the original paired $SU(2)$ color/gamma/Fock data, the actual physical
Hilbert subspace is nonzero:
\begin{equation}\label{eq:physical-nonzero}
 \mathscr H_{\mathrm{phys}}\neq\{0\}.
\end{equation}
Indeed, a source-defined gauge-invariant fermionic vacuum multiplied by
a compactly supported radial profile defines a nonzero member.
\end{theorem}
\begin{proof}[Mathematical construction]
Let $v_0\in\mathcal F$ be the nonzero original gauge-fixed fermionic
vacuum obtained in the source construction, and choose a continuous radial
compactly supported scalar profile $\chi(\|x\|^2)$ with $\chi(0)=1$.
The concrete source takes a smooth-transition function of
$2-2\|x\|^2$. Define
\begin{equation}\label{eq:radial-bump}
 \psi_0(x):=\chi(\|x\|^2)v_0.
\end{equation}
Compact support and continuity imply $\psi_0\in L^2$. The color action
preserves $\|x\|$, while the original fermionic unitary fixes $v_0$;
thus $\psi_0$ obeys the physical covariance condition. Finally,
$\psi_0(0)=v_0\ne0$, and by continuity $\psi_0$ is nonzero on an open
neighborhood of the origin of positive volume. Therefore its $L^2$ class
is nonzero. The last step is essential: nonzero value at \emph{one}
point alone is not evidence that an $L^2$ equivalence class is nonzero.
\end{proof}
\formalrecord{Original G4--K2B declarations
\texttt{radialBumpL2\_ne\_zero},
\texttt{radialBumpL2\_mem\_physicalSpace},
\texttt{pairedPhysicalSpace\_ne\_bot}; historical FCP strict
\href{https://github.com/etblink/Foundational-Convergence-Program/actions/runs/37894319431}{\#37894319431}
(8 October 2026). Original source port is in
\path{OperatorDependencies/SU2BFSSG4K2BNonzeroPhysicalL2Probe.lean}.
Atlas G-02.}

In November, the unified integration program also qualified a nonzero
\emph{finite original Fock-vector} invariant and a nonzero original
fermionic gauge projection. This provides an algebraically useful
independent construction, but its finite fiber fixed-space conclusion
is not the same theorem as \eqref{eq:physical-nonzero} about the full
$L^2$ physical wavefunction space.
\formalrecord{\path{OP2cNonzeroOriginalGaugeVacuum.lean}; integrated strict
\href{https://github.com/etblink/math/actions/runs/42530241914}{\#42530241914}.
Atlas G-03.}

\section{The Gauss-law kernel is more than a formal slogan}
The traditional Gauss constraint suggests that physical states are those
annihilated by the three infinitesimal color generators. On unbounded
operators, this statement needs domains. The historical FCP work
constructed the actual \emph{maximal strong derivative} generators of
the original full-$L^2$ group action, then established the relation of
their joint kernel to the global invariant space.

\begin{sourceaccount}[Gauss-kernel/physical-space equivalence;
historically qualified]
For the original paired SU(2) action, an original full-$L^2$ state is
physical if and only if it lies in the domains of all three actual color
generators and is annihilated by each generator.
\end{sourceaccount}
\begin{sourceexplanation}
The forward direction follows from the fact that the orbit of an
invariant vector is constant, so its strong derivative along every
one-parameter color subgroup exists and vanishes. Conversely, the
historically qualified one-parameter generator theorem identifies
kernel vectors with vectors fixed by each of the corresponding flows.
A separate \emph{exact} SU(2) Euler decomposition establishes that those
flows generate the whole genuine matrix group, not merely a dense
subgroup. Consequently joint zero-generator states are globally fixed
and satisfy \eqref{eq:physical-condition}. This is a characterization
of a joint kernel, \emph{not} equality between maximal generator domains
and $\mathscr D_{\mathrm{phys}}$.
\end{sourceexplanation}
\formalrecord{FCP historical Gauss source theorem family
\texttt{L6H5E8AA}, strict
\href{https://github.com/etblink/Foundational-Convergence-Program/actions/runs/38994451291}{\#38994451291}
(16 October 2026); Euler decomposition \texttt{H5E8Y} strict
\href{https://github.com/etblink/Foundational-Convergence-Program/actions/runs/38992065735}{\#38992065735}.
The chronological source handoff is \path{FCP_BFSS_SU2_GAUSS_KERNEL_COMPLETION_AND_SPECTRAL_RETURN_2026_10_16.md}.}

\section{Reader's checkpoint}
We now know which vectors are physical, that the physical space is
closed and nonzero, and that the original full-$L^2$ Gauss law has an
accepted strong-generator characterization. None proves that the
physical charge domain is dense, nor that a formal Hamiltonian preserves
a Spin(9) isotypic sector. Those questions demand a separate
closed-operator construction.
